% Standalone section for problems/3.2/proof.tex — \input this file.
% Written 2026-07-31 (life window, Fable campaign).  Every numerical claim has a
% runnable audit in problems/3.2/research/scripts/ (named inline).
% NOTE: NEW file (single-writer discipline); do not merge into proof.tex by editing.

\section{Value distribution of the Ap\'ery sequence in one period:
unconditional multiplicity, image and energy bounds}
\label{sec:value-distribution}

Throughout this section $p\ge5$ is prime, $b_n=\sum_k\binom nk^2\binom{n+k}k^2$,
$N_p(c)=\#\{0\le r\le p-1:\ b_r\equiv c\ (p)\}$, and
$E(p)=\sum_{c\in\mathbb F_p}N_p(c)^2$ is the one-period collision energy
(empirically $E(p)=3p+O(\sqrt p)$; script \texttt{fable\_vertical\_value\_law.py}).
The engine of this section is the \emph{non-autonomy} of the Ap\'ery recurrence
$(m+1)^3b_{m+1}=P(m)b_m-m^3b_{m-1}$, $P(m)=34m^3+51m^2+27m+5$: the moving
coefficient acts as the derivative-substitute that Weil/Stepanov technology
lacks here (cf.\ \S\ref{sec:marked-scalar-structure} for why the bounded-conductor
routes are closed).  Transfer notation: $b_{r+d}=A_d(r)b_r+B_d(r)b_{r-1}$ with
$A_d,B_d\in\mathbb F_p(x)$, poles confined to $\{-1,\dots,-d\}$ of order $\le3$.

\subsection{Statements}

\begin{theorem}[image]\label{thm:image}
$\#\{b_r\bmod p:\ 0\le r<p\}\ \ge\ ((p-1)/4)^{1/3}$ for every prime $p$.
\end{theorem}

\begin{theorem}[multiplicity]\label{thm:mult}
There is an absolute $C_0$ with $\max_{c\ne0}N_p(c)\le C_0\,p/\log p$ for every prime $p$.
\end{theorem}

\begin{theorem}[energy]\label{thm:energy}
$E(p)\ \ll\ p^2/\log p$ for every prime $p$.  Via Parseval
($\sum_{h\ne0}|C_p(h)|^2=pE(p)-p^2$) this gives the average vertical bound
$\bigl(\tfrac1{p-1}\sum_{h\ne0}|C_p(h)|^2\bigr)^{1/2}\ll p/\sqrt{\log p}$.
\end{theorem}

\begin{theorem}[$p^{3/4}$ off a thin prime set]\label{thm:threequarters}
There is a set $\mathfrak E$ of primes with
$\#\{q\in\mathfrak E: q\le X\}=O(X^{3/4+o(1)})$ such that for $p\notin\mathfrak E$
and all $c\ne0$: $N_p(c)\ll p^{3/4}(\log p)^{O(1)}$ and
$E(p)\ll p^{7/4}(\log p)^{O(1)}$.
\end{theorem}

\begin{theorem}[short lags]\label{thm:shortlag}
For $\Delta\le c_0\log p$ (all $p$; and for $\Delta\le p^{1/3}$ when
$p\notin\mathfrak E$), $\sum_{1\le d\le\Delta}\#\{r:\ b_{r+d}\equiv b_r\}
\ \le\ 2p+6\sqrt{p\,\Delta^3}$.
\end{theorem}

\begin{theorem}[universal nondegeneracy, (UN)]\label{thm:UN}
For every prime $p\ge5$ and all $1\le h<k<p$ the rational function
$\Psi_{h,k}=(1-A_h)B_k-(1-A_k)B_h$ is not identically zero over $\mathbb F_p$;
its reduced numerator has degree exactly $3k$ (measured; claimed $\le3(k-1)$ up
to the normalizing unit).
\end{theorem}

\noindent\emph{Proof status.}  Complete 17-lemma proof delivered (session round
R4b, archived \texttt{chatgpt-answers/Q6311-94300c1e.md}; full display formulas
on the Notion source page — transcription into this section is a pending
editorial task).  Chain: cleared continuants $Q_m$; $\Psi=\mathrm{unit}\cdot
\Delta_{h,k}$ with $\Delta$ polynomial; $Q_m\not\equiv0$ for $m<p$; continuant
reflection via tridiagonal $J$-conjugation; Pl\"ucker identity and subtractive
Euclidean descent $(h,k)\to(\min(h,k-h)\dots)$; terminal central-return case
$(d,2d)$, $C^d=\pm I$, excluded by a power-series expansion at infinity with
leading coefficients $5d/64$ ($\epsilon=-1$, at $z^2$) and $-75d^3/1024$
($\epsilon=+1$, at $z^4$), nonzero for $p\ge7$; $p=5$ explicit.  Machine
verification: the terminal expansion checked for ALL 277 applicable $(p,d,
\epsilon)$ with $7\le p<500$ — zero failures (\texttt{un\_proof\_check.py});
the statement itself swept clean over $\approx10^6$ cases
(\texttt{universal\_nonvanish.py}, \texttt{stress\_test.py}).

\begin{theorem}[all-prime $p^{3/4}$ and $p^{5/3}$]\label{thm:final}
For every prime $p\ge5$: (i) $N_p(c)\le 8\,p^{3/4}$ for every $c\in\mathbb F_p$
(including $c=0$), by the block/H\"older triple count driven by
Theorem~\ref{thm:UN}; (ii) $E(p)\ll p^{5/3}$, by the $c$-uniform
triple-convexity count ($\Sigma_c m_c^3\ll p^2H$ at block length $H$,
Cauchy--Schwarz, $H=p^{1/3}$); (iii) hence
$\bigl(\tfrac1{p-1}\sum_{h\ne0}|C_p(h)|^2\bigr)^{1/2}\ll p^{5/6}$.
Theorems~\ref{thm:mult} and~\ref{thm:threequarters} are superseded by (i)
(kept above for their elementary self-contained proofs).
\end{theorem}

\begin{proposition}[methodological ceiling]\label{prop:ceiling}
An abstract block-permutation adversary realizes every consequence of the
one-base certificate method (value-independent nonzero certificates of degree
$O(k\log p)$ per lag pair, reflection, spacing) with $E\asymp p^{5/3}$.  Hence
no refinement of the present architecture alone reaches $E\ll p^{3/2}$; the
$\rho$-ladder (family-average root count $R_H\ll H^\rho$: $\rho=3$ now,
$\rho=2\Rightarrow p^{3/2}$, $\rho\to1\Rightarrow p^{1+o(1)}$) requires
filtering candidate kernels to the distinguished Ap\'ery orbit (two-base
compatibility, anchored backward transfer, or spectral input).  A Kummer-order
stratification experiment (40 primes to $1.8\times10^5$) shows no bounded-order
concentration of collision lags — the Mellin/Kummer route has no self-twist
shortcut, and the sub-$5/3$ wall is structural for the present methods.
\end{proposition}

\begin{proposition}[downstream limitation — do not overclaim]\label{prop:limitation}
Single-prime vertical bounds enter the gcd problem of
Problem~3.2 only through $|Z_p|\le E(p)^{1/2}$.  Since the zero-fibre bound
$|Z_p|\le3p^{2/3}$ is already stronger than $E(p)^{1/2}$ for every exponent
proved here, \emph{Theorems~\ref{thm:mult}--\ref{thm:shortlag} yield no
improvement of the density-one theorem nor of the pointwise problem}.  Moreover
no family of per-prime vertical bounds can settle the pointwise problem: the
singleton-alignment countermodel ($Z_p=\{n_0-p\}$) satisfies every vertical
constraint yet gives $H(n_0)\gg n_0/\log n_0$.  The remaining interface is
cross-prime: with $|Z_p|\ll p^\alpha$, a $k$-th cross-prime moment bound at the
independence scale with $k>1/(1-\alpha)$ suffices ($\alpha=2/3$: fourth moment).
\end{proposition}

\subsection{The collision mechanism}

Fix $c\ne0$ and $1\le r<r+d\le p-2$ with $b_r\equiv b_{r+d}\equiv c$.  From
$c=A_d(r)c+B_d(r)b_{r-1}$:
\emph{Type I} ($B_d(r)\ne0$): $b_{r-1}=c\,\sigma_d(r)$ with
$\sigma_d=(1-A_d)/B_d$ — the predecessor is determined.
\emph{Type II} ($B_d(r)=0$): $A_d(r)=1$, so $r$ roots the numerator of $B_d$
(degree $\le4d$).  If $r$ has two type-I lags $d_1<d_2$ then
$\sigma_{d_1}(r)=\sigma_{d_2}(r)$, i.e.\ $r$ roots the numerator of
$\Psi_{d_1,d_2}$, of degree $\le9(d_1+d_2)$.
Machine check: all $60$ multi-lag fiber elements at $p=101,199$ root the
predicted polynomial (\texttt{psi\_e2e.py}); at $p=101$ the $49$ type-II
collision pairs are exactly the reflection pairs $r\leftrightarrow p-1-r$.

\subsection{Polar constants and nondegeneracy}

\begin{lemma}[rank-one polar identity]\label{lem:polar}
$(-5,\,1)\cdot T_{d-1}(-d)=d^3\,(-b_d,\;b_{d-1})$.  Consequently at $x=-d$ the
transfer entries have poles of order exactly $3$ (over $\mathbb Q$) with
\[
c_d:=\lim_{x\to-d}(x+d)^3B_d(x)=d^3\,b_{d-1},\qquad
a_d:=\lim_{x\to-d}(x+d)^3A_d(x)=-d^3\,b_d .
\]
\end{lemma}

\begin{proof}
Split $M(x)=(x+1)^{-3}(1,0)^{\mathsf T}(P(x),-x^3)+\left(\begin{smallmatrix}0&0\\1&0\end{smallmatrix}\right)$;
$T_{d-1}$ is regular at $-d$, so the polar part of $T_d=M(x+d-1)T_{d-1}$ at
$-d$ is $(1,0)^{\mathsf T}(-5,1)T_{d-1}(-d)$.  Induction using
$P(-1-n)=-P(n)$ and the recurrence: base
$(-5,1)M(-2)=(-584,40)=2^3(-b_2,b_1)$; step
$d^3(-b_d,b_{d-1})M(-d-1)=(d+1)^3(-b_{d+1},b_d)$ is exactly the recurrence.
(Exact-arithmetic verification $d\le30$: \texttt{cd\_test.py}, \texttt{ad\_test.py}.)
\end{proof}

\begin{lemma}[two endpoint witnesses]\label{lem:witness}
Let $1\le d_1<d_2$, $\delta=d_2-d_1$.  The $(x+d_1)^{-3}$-coefficient of
$\Psi_{d_1,d_2}$ equals $(b_\delta-1)\,d_1^3\,b_{d_1-1}$, and the
$(x+d_2)^{-3}$-coefficient equals $d_2^3\,(b_{d_2-1}-b_{\delta-1})$.
Both are nonzero over $\mathbb Q$ ($b$ is strictly increasing with $b_1=5$).
Hence $\Psi_{d_1,d_2}\not\equiv0$ over $\mathbb Q$ for every pair, and over
$\mathbb F_p$ whenever $p$ fails to divide the corresponding integer.
\end{lemma}

\begin{proof}
First coefficient: $\Psi=(B_{d_2}-B_{d_1})-K$, $K=\det T_{d_1}\cdot
B_\delta(x+d_1)$ is regular at $-d_1$ (pole$^3\times$zero$^3$;
$B_\delta\in x^3\mathbb F_p[[x]]$ always by $B_k(x)=-x^3(x+1)^{-3}A_{k-1}(x+1)$
— over $\mathbb F_p$ the cubic coefficient may vanish, membership in the cube
ideal suffices), and $B_{d_2}\sim A_\delta(0)B_{d_1}=b_\delta B_{d_1}$ there.
Second coefficient: the reflected scalar identity (negative-window transfer
reflected by $P(-1-n)=-P(n)$) evaluates
$(1-A_{d_1}(-d_2))b_{d_2-1}-b_{d_2}B_{d_1}(-d_2)=b_{d_2-1}-b_{\delta-1}$;
verified for $145$ pairs at $p=101$ including direct symbolic Laurent limits
(\texttt{K\_test.py} and the audited snippet in \texttt{Q6306.md}).
\end{proof}

\subsection{Proofs of Theorems \ref{thm:image}--\ref{thm:shortlag}}

\emph{Theorem \ref{thm:image}.}  If the image has size $D$, some state
$(b_r,b_{r-1})=(x,y)$ repeats $\ge(p-1)/D^2$ times; for those $r$,
$b_{r+1}=(P(r)x-r^3y)(r+1)^{-3}$ is a fixed rational function of degree
$\le4$, nonconstant unless $(x,y)=(0,0)$ (excluded: no two consecutive zeros,
by backward propagation to $b_0=1$), taking each value $\le4$ times inside the
image: $(p-1)/D^2\le4D$. \qed

\emph{Theorem \ref{thm:mult}.}  Order the fiber; $d_1(r)<d_2(r)$ the next-two
lags.  $\sum_r(d_1+d_2)(r)\le3p$, so $\ge N/2-O(1)$ elements have
$d_1+d_2\le D':=12p/N$ (the earlier $3N/4$ claim is false — adversarial
near-uniform fibers; $N/2$ suffices).  If $N\ge C_0p/\log p$ then
$3.526\,k\le0.85\log p$ for every index $k\le D'$ in play, so the integers
$b_{d_1-1}$ and $b_\delta-1$ lie in $(0,p)$ ($b_k<(1+\sqrt2)^{4k}$) and every
witness of Lemma~\ref{lem:witness} is a unit mod $p$: every pair in play is
nondegenerate.  Capacity: type-II $\sum_{d\le D'}O(d)=O(D'^2)$ plus type-I
$\sum_{d_1<d_2\le D'}9(d_1+d_2)\le5D'^3\ll(\log p)^3<N/4$.  Contradiction. \qed

\emph{Theorem \ref{thm:energy}.}  $E\le N_p(0)^2+\max_{c\ne0}N_p(c)\cdot p$
with $N_p(0)\le3p^{2/3}$. \qed

\emph{Theorem \ref{thm:threequarters}.}  Run the same count at
$D'=12p/N$, $N\ge p^{3/4+\varepsilon}$.  A pair escapes both witnesses of
Lemma~\ref{lem:witness} only if $p\mid\gcd(b_\delta-1,\ b_{d_2-1}-b_{\delta-1})$
(or the $b_{d_1-1}$-analogue) — fixed nonzero integers of log-size $O(d_2)$.
$\#\{\text{bad }q\le X\}\le\sum_{d_1<d_2\le X^{1/4}}\omega(\gcd)=O(X^{3/4+o(1)})$.
For $p\notin\mathfrak E$ no escape exists and the counting is verbatim.
(Sweep $d_1<20$, $d_2\le50$: every prime factor of every such gcd is
$\le89\approx2.4\,d_2$ — the visible exceptional set is empty;
\texttt{G\_wide\_sweep.py}, \texttt{bracket\_test.py}.) \qed

\emph{Theorem \ref{thm:shortlag}.}  The $\Psi$ polynomials are independent of
$c$, so $\sum_r\binom{L_\Delta(r)}2\le\sum_{d_1<d_2\le\Delta}9(d_1+d_2)\le
5\Delta^3$ where $L_\Delta(r)$ counts same-fiber lags $\le\Delta$;
Cauchy--Schwarz on $T_\Delta=\sum_rL_\Delta(r)$ gives
$T_\Delta\le2p+6\sqrt{p\Delta^3}$. \qed

\subsection{Provenance and audit trail}

All identities were discovered and verified in exact arithmetic before use
(scripts named above; see \texttt{FABLE\_NOTES\_energy\_bootstrap.md} for the
session log).  Negative literature search (audited): no prior maximum-fiber,
image, or energy bound of this type for Ap\'ery or general $P$-recursive
sequences mod $p$ was located
(closest: Garaev--Luca--Shparlinski, JCTA 113 (2006), residue classes /
additive bases; Luca--Shparlinski, JLMS 78 (2008), prime factors of $b_n$;
Ostafe--Shparlinski, Canad.\ J.\ Math.\ 74 (2022), $C$-finite families).
We therefore state novelty as ``we are not aware of prior bounds of this
form''.
